% GENERATED FILE. Edit canonical lessons, then run npm run generate:books.
% canonical-source-hash: fnv1a64:e14073e3b4fa282e
% canonical-lessons: TA-C439-vantu, TA-C439-paci, TA-C439-marakkanru, TA-C439-mulai, TA-C439-acai

\chapter{Sprouts and Moss}
\label{ch:ta-sprouts-and-moss}
\hlchaptermodality{439}

\begin{chapteropening}
\textbf{By the end of this chapter:} \emph{I can talk about beetles, moss and saplings, and seeds sprouting and leaves stirring.}

Complete the last lesson of chapter 439: I can talk about beetles, moss and saplings, and seeds sprouting and leaves stirring.
\end{chapteropening}

\section[\texorpdfstring{\ta{வண்டு} (vaṇṭu)}{வண்டு (vaṇṭu)}]{\ta{வண்டு} (vaṇṭu) — a beetle; a bee}

\label{lesson:TA-C439-vantu}



\begin{quote}
Before the new one: say the Tamil for a lizard, then the Tamil for a scorpion.
\end{quote}

\subsection*{You'll want to know: \ta{வண்டு}}
\textbf{\ta{வண்டு}} — \emph{vaṇṭu} — \textquotedblleft{}a beetle; a bee\textquotedblright{}.

\begin{quote}
\ta{பூக்களைச்} \ta{சுற்றி} \ta{வண்டுகள்} \ta{பறந்தன}.
\end{quote}

\begin{quote}
\emph{pūkkaḷaic cuṟṟi vaṇṭukaḷ paṟantaṉa.}
\end{quote}

\begin{quote}
Beetles flew around the flowers.
\end{quote}

Any buzzing insect that visits flowers.

\begin{cousinweb}
\textbf{\ta{வண்டு}} (\emph{vaṇṭu}) is native Dravidian (DEDR 5239), \textquotedblleft{}a bee, a beetle\textquotedblright{}.
\end{cousinweb}

\begin{grammarlens}[title={What you've built}]
One more word for getting around.
\end{grammarlens}

\subsection*{Your turn}
\begin{itemize}
\raggedright
  \item \emph{Say it:} \emph{vaṇṭu}
  \item \emph{Say it:} \emph{vaṇṭu}, once more
  \item \emph{Say it:} \emph{pūkkaḷaic cuṟṟi vaṇṭukaḷ paṟantaṉa.} — Beetles flew around the flowers.
  \item \emph{Recall:} say the Tamil for a lizard, then the Tamil for a scorpion, then say \emph{vaṇṭu} again
\end{itemize}

\begin{tcolorbox}[breakable,colback=teal!4,colframe=teal!35!black,title={Before you move on}]
What does \ta{வண்டு} mean? (\textquotedblleft{}A beetle; a bee\textquotedblright{}.) Say it once more.
\end{tcolorbox}
\section[\texorpdfstring{\ta{பாசி} (pāci)}{பாசி (pāci)}]{\ta{பாசி} (pāci) — moss; algae}

\label{lesson:TA-C439-paci}



\begin{quote}
Before the new one: say the Tamil for a scorpion, then the Tamil for a beetle.
\end{quote}

\subsection*{You'll want to know: \ta{பாசி}}
\textbf{\ta{பாசி}} — \emph{pāci} — \textquotedblleft{}moss; algae\textquotedblright{}.

\begin{quote}
\ta{குளத்தின்} \ta{படிகளில்} \ta{பாசி} \ta{இருந்ததால்} \ta{வழுக்கியது}.
\end{quote}

\begin{quote}
\emph{kuḷattiṉ paṭikaḷil pāci iruntatāl vaḻukkiyatu.}
\end{quote}

\begin{quote}
It was slippery because there was moss on the steps of the pond.
\end{quote}

From the same root as \ta{பச்சை}, green.

\begin{cousinweb}
\textbf{\ta{பாசி}} (\emph{pāci}) is native Dravidian (DEDR 3821), \textquotedblleft{}that which is green; moss, lichen, seaweed\textquotedblright{}.
\end{cousinweb}

\begin{grammarlens}[title={What you've built}]
One more word for getting around.
\end{grammarlens}

\subsection*{Your turn}
\begin{itemize}
\raggedright
  \item \emph{Say it:} \emph{pāci}
  \item \emph{Say it:} \emph{pāci}, once more
  \item \emph{Say it:} \emph{kuḷattiṉ paṭikaḷil pāci iruntatāl vaḻukkiyatu.} — It was slippery because there was moss on the steps of the pond.
  \item \emph{Recall:} say the Tamil for a scorpion, then the Tamil for a beetle, then say \emph{pāci} again
\end{itemize}

\begin{tcolorbox}[breakable,colback=teal!4,colframe=teal!35!black,title={Before you move on}]
What does \ta{பாசி} mean? (\textquotedblleft{}Moss; algae\textquotedblright{}.) Say it once more.
\end{tcolorbox}
\section[\texorpdfstring{\ta{மரக்கன்று} (marakkaṉṟu)}{மரக்கன்று (marakkaṉṟu)}]{\ta{மரக்கன்று} (marakkaṉṟu) — a sapling}

\label{lesson:TA-C439-marakkanru}



\begin{quote}
Before the new one: say the Tamil for a beetle, then the Tamil for moss.
\end{quote}

\subsection*{You'll want to know: \ta{மரக்கன்று}}
\textbf{\ta{மரக்கன்று}} — \emph{marakkaṉṟu} — \textquotedblleft{}a sapling\textquotedblright{}.

\begin{quote}
\ta{பள்ளி} \ta{மாணவர்கள்} \ta{நூறு} \ta{மரக்கன்றுகளை} \ta{நட்டார்கள்}.
\end{quote}

\begin{quote}
\emph{paḷḷi māṇavarkaḷ nūṟu marakkaṉṟukaḷai naṭṭārkaḷ.}
\end{quote}

\begin{quote}
The schoolchildren planted a hundred saplings.
\end{quote}

\ta{மரம்}, a tree, + \ta{கன்று}, a young one: \ta{மரம்} drops -\ta{ம்} and doubles \ta{க்}.

\begin{cousinweb}
\textbf{\ta{மரக்கன்று}} (\emph{marakkaṉṟu}) is built from \ta{மரம்}, \textquotedblleft{}a tree\textquotedblright{} (DEDR 4711), and \ta{கன்று}, \textquotedblleft{}a calf, a sapling, a young tree\textquotedblright{} (DEDR 1411).
\end{cousinweb}

\begin{grammarlens}[title={What you've built}]
One more word for getting around.
\end{grammarlens}

\subsection*{Your turn}
\begin{itemize}
\raggedright
  \item \emph{Say it:} \emph{marakkaṉṟu}
  \item \emph{Say it:} \emph{marakkaṉṟu}, once more
  \item \emph{Say it:} \emph{paḷḷi māṇavarkaḷ nūṟu marakkaṉṟukaḷai naṭṭārkaḷ.} — The schoolchildren planted a hundred saplings.
  \item \emph{Recall:} say the Tamil for a beetle, then the Tamil for moss, then say \emph{marakkaṉṟu} again
\end{itemize}

\begin{tcolorbox}[breakable,colback=teal!4,colframe=teal!35!black,title={Before you move on}]
What does \ta{மரக்கன்று} mean? (\textquotedblleft{}A sapling\textquotedblright{}.) Say it once more.
\end{tcolorbox}
\section[\texorpdfstring{\ta{முளை} (muḷai)}{முளை (muḷai)}]{\ta{முளை} (muḷai) — to sprout}

\label{lesson:TA-C439-mulai}



\begin{quote}
Before the new one: say the Tamil for moss, then the Tamil for a sapling.
\end{quote}

\subsection*{You'll want to know: \ta{முளை}}
\textbf{\ta{முளை}} — \emph{muḷai} — \textquotedblleft{}to sprout\textquotedblright{}.

\begin{quote}
\ta{மழைக்குப்} \ta{பிறகு} \ta{விதைகள்} \ta{முளைத்தன}.
\end{quote}

\begin{quote}
\emph{maḻaikkup piṟaku vitaikaḷ muḷaittaṉa.}
\end{quote}

\begin{quote}
After the rain the seeds sprouted.
\end{quote}

Past (it …): \ta{முளைத்தது}, \emph{muḷaittatu}.

\ta{முளைகட்டிய} \ta{பயறு} is sprouted gram.

\begin{cousinweb}
\textbf{\ta{முளை}} (\emph{muḷai}) is native Dravidian (DEDR 4997), \textquotedblleft{}to spring up, to sprout\textquotedblright{}.
\end{cousinweb}

\begin{grammarlens}[title={What you've built}]
One more word for getting around.
\end{grammarlens}

\subsection*{Your turn}
\begin{itemize}
\raggedright
  \item \emph{Say it:} \emph{muḷai}
  \item \emph{Say it:} \emph{muḷai}, once more
  \item \emph{Say it:} \emph{maḻaikkup piṟaku vitaikaḷ muḷaittaṉa.} — After the rain the seeds sprouted.
  \item \emph{Recall:} say the Tamil for moss, then the Tamil for a sapling, then say \emph{muḷai} again
\end{itemize}

\begin{tcolorbox}[breakable,colback=teal!4,colframe=teal!35!black,title={Before you move on}]
What does \ta{முளை} mean? (\textquotedblleft{}To sprout\textquotedblright{}.) Say it once more.
\end{tcolorbox}
\section[\texorpdfstring{\ta{அசை} (acai)}{அசை (acai)}]{\ta{அசை} (acai) — to sway; to stir}

\label{lesson:TA-C439-acai}



\begin{quote}
Before the new one: say the Tamil for a sapling, then the Tamil for to sprout.
\end{quote}

\subsection*{You'll want to know: \ta{அசை}}
\textbf{\ta{அசை}} — \emph{acai} — \textquotedblleft{}to sway; to stir\textquotedblright{}.

\begin{quote}
\ta{காற்றில்} \ta{இலைகள்} \ta{மெதுவாக} \ta{அசைந்தன}.
\end{quote}

\begin{quote}
\emph{kāṟṟil ilaikaḷ metuvāka acaintaṉa.}
\end{quote}

\begin{quote}
The leaves swayed slowly in the wind.
\end{quote}

Past (it …): \ta{அசைந்தது}, \emph{acaintatu}.

\ta{அசைந்தது} is it moved; \ta{அசைத்தேன்}, I moved something.

\begin{cousinweb}
\textbf{\ta{அசை}} (\emph{acai}) is native Dravidian (DEDR 37), \textquotedblleft{}to move, to stir; to shake\textquotedblright{}.
\end{cousinweb}

\begin{grammarlens}[title={What you've built}]
One more word for getting around.
\end{grammarlens}

\subsection*{Your turn}
\begin{itemize}
\raggedright
  \item \emph{Say it:} \emph{acai}
  \item \emph{Say it:} \emph{acai}, once more
  \item \emph{Say it:} \emph{kāṟṟil ilaikaḷ metuvāka acaintaṉa.} — The leaves swayed slowly in the wind.
  \item \emph{Recall:} say the Tamil for a sapling, then the Tamil for to sprout, then say \emph{acai} again
\end{itemize}

\begin{tcolorbox}[breakable,colback=teal!4,colframe=teal!35!black,title={Before you move on}]
What does \ta{அசை} mean? (\textquotedblleft{}To sway; to stir\textquotedblright{}.) Say it once more.
\end{tcolorbox}
